\chapter{Introduction}
\label{ch:introduction}

\section{Background}

The introduction should establish the problem, its significance, and the context
needed to understand the thesis. Begin with the broader research area, narrow the
discussion to the specific problem, and define the scope of the investigation.
Avoid turning this section into a second literature review; reserve detailed
comparison of prior methods for \Cref{ch:literature-review}.

Factual claims should be supported by appropriate sources. This sentence
demonstrates a compressed numeric citation \parencite{example2025,example2024};
citations appear as bracketed numbers such as [1] in the compiled thesis.

\section{Aim and Objectives}

State the overall aim in one precise sentence, then translate it into objectives
that can be evaluated by the methods and results. Each objective should correspond
to a later methodological step and a reported finding.

An objectives list can be used when it improves readability:

\begin{enumerate}[label=\textbf{O\arabic*:},leftmargin=*]
  \item define a reproducible problem setting;
  \item compare representative methods under consistent conditions; and
  \item analyse uncertainty, limitations, and generalisability.
\end{enumerate}

\subsection{Research Questions}

Research questions should be specific, answerable, and aligned with the objectives.
If hypotheses are used, state them after the relevant question and define the
evidence that would support or challenge each one.

\subsubsection{Scope and Terminology}

Define the population, data, task, and evaluation boundary that apply throughout
the thesis. Introduce specialised terms once and use them consistently; a glossary
or abbreviation list is preferable when many recurring terms are required.

Technical terms may be clarified in a footnote.\footnote{Footnotes should be concise and should not be used as a substitute for essential argumentation or citation.}

\section{Thesis Organisation}

\Cref{ch:literature-review} reviews the relevant literature. \Cref{ch:study-one}
demonstrates a self-contained study with its own introduction, methods, results,
and discussion. \Cref{ch:conclusions} integrates the conclusions and identifies
directions for future research.
